\section{Proofs for decomposition certificates}
\label{app:decomp}

\subsection{Configurations and shell partitions}
\label{app:decomp:existence}

\begin{lemma}[Unfolding the local inequalities]\label{decomp:unfolding}
Use \eqref{decomp:common-slope-child}--\eqref{decomp:intercepts}.
A configuration chooses a local leaf $B_t$ and a point $z^t\in B_t$ for
each bag. For each nonroot bag it also chooses a cell $E_t$ such that
$z^t_{S_t}\in E_t$ and $E_t\cap(B_{p(t)})_{S_t}\ne\varnothing$.
Then
\begin{equation}\label{decomp:configuration-value}
 \ell_r=\min_{\text{configurations}}
 \left\{\sum_t\sum_{c\text{ assigned to }t}f_{c,(B_t)_c}(z^t_c)
 +\sum_{t\ne r}\lambda_t^{\mathsf T}
                    (z^{p(t)}_{S_t}-z^t_{S_t})\right\}.
\end{equation}
\end{lemma}
\begin{proof}
For a fixed root leaf and point, each child term selects a cell meeting
the root leaf projection and pays its intercept. The intercept is the
minimum over the child's leaf and point restricted to that cell.
Substitute its definition. The child's positive slope term, evaluated at
the parent point, combines with the negative slope term in its intercept
to give the corresponding summand in
\eqref{decomp:configuration-value}. Repeat for every child.
At each step, child choices have independent coordinate copies; no
equality of copies is imposed. A sum of independent minima is the minimum
over the product of their choices. There are finitely many boxes and all
point domains are compact. Continuity supplies every minimum and completes
the recursive substitution.
\end{proof}

For a box $B$, write $\omega(B)$ for its largest side length. For a point
$p$, let $Q(p,a)=p+[-a,a]^d$.
For an empty coordinate set, the coordinate box is the singleton in
$\R^0$, its partition has one cell, and $\omega(B)$ and all
empty-coordinate norms are zero. Interiors are always taken relative
to the corresponding coordinate space. These conventions cover empty
separators.

\begin{lemma}[Shell partition]\label{decomp:shells}
Let $X$ be a nondegenerate box in $\R^d$ with side lengths at most $s_0>0$,
$p\in X$, $h>0$, and $\theta=2^{-\mu}$ for an integer $\mu\ge0$.
There is a rectangular partition of $X$ with at most
$(J+1)(4/\theta)^d$ members, where
$J=\max\{0,\lceil\log_2(s_0/h)\rceil\}$, such that every member satisfies
\begin{equation}\label{decomp:shell-width}
 \omega(B)\le\max\{h,\theta\dist_\infty(B,p)\}
              \le h+\theta\dist_\infty(B,p).
\end{equation}
\end{lemma}
\begin{proof}
Partition $Q(p,h)$ into $2^d$ cubes of side $h$ with vertex $p$.
For $j=1,\ldots,J$, put a grid of side
$g_j=\theta2^{j-1}h$ on $Q(p,2^jh)$, anchored at its lower corner.
There are $4/\theta$ grid cells per coordinate. The boundary of
$Q(p,2^{j-1}h)$ is on grid lines, because $1/\theta$ is an integer.
Retain the grid cells outside the interior of that smaller cube, omitting
the cells contained in it. Each retained cell has
$\dist_\infty(B,p)\ge2^{j-1}h$ and side $g_j$, so it satisfies
\eqref{decomp:shell-width}.

These levels cover $Q(p,2^Jh)$ with disjoint interiors. This cube contains
$X$, because $p\in X$ and $2^Jh\ge s_0$. Intersect the cells with $X$ and
discard degenerate intersections. Clipping decreases widths and increases
distance from $p$. The nondegenerate intersections still cover $X$:
they cover its interior, and their finite closed union therefore contains
its closure. Each level contributes at most $(4/\theta)^d$ boxes.
The central level contributes $2^d\le(4/\theta)^d$.
\end{proof}

\begin{proof}[Proof of Theorem~\ref{decomp:small-certificate}]
Use Lemma~\ref{decomp:shells} in every bag and separator, centered at the
corresponding coordinates of $x^*$. Use the slopes
\eqref{decomp:gradient-slopes} and the maximal intercepts
\eqref{decomp:intercepts}. We bound every configuration in
\eqref{decomp:configuration-value}.

For each coordinate $i$, let $\operatorname{top}(i)$ be the bag nearest
the root that contains it. Define a consistent point $x\in D$ by
$x_i=z^{\operatorname{top}(i)}_i$. Connectedness of the bags containing
$i$ ensures that its other copies are connected to this copy through
at most $k-1$ edges containing $i$ in their separators. Put
\[
 X_*=\|x-x^*\|_2,\quad
 \rho_t=\|x_{V_t}-x^*_{V_t}\|_\infty,\quad
 r_t=\|x_{V_t}-x^*_{V_t}\|_2,
\]
\[
 d_t=\|z^t_{S_t}-z^{p(t)}_{S_t}\|_\infty\quad(t\ne r),\quad d_r=0,
 \qquad \delta_t=\|z^t-x_{V_t}\|_\infty.
\]
Empty-separator drifts are zero. Since every variable occurs in at most
$k$ bags,
\begin{equation}\label{decomp:bag-radii}
 \sum_t\rho_t^2\le\sum_tr_t^2\le kX_*^2.
\end{equation}

Let $\operatorname{anc}_j(t)$ denote the $j$th ancestor, with missing
ancestor drifts set to zero. Telescoping each coordinate's copies gives
\begin{equation}\label{decomp:bag-drift}
 \delta_t\le\sum_{j=0}^{k-2}d_{\operatorname{anc}_j(t)}.
\end{equation}
The nonseparator coordinates have their topmost bag at $t$, so only
separator coordinates can differ between $z^t$ and $x_{V_t}$.

The width estimate in \eqref{decomp:shell-width}, evaluated at the bag
copy, gives
\[
 \omega(B_t),\ \omega(E_t)
          \le2h+\theta(\rho_t+\delta_t),
\]
where the harmless $2h$ also covers a zero-dimensional separator.
Choose a point in $E_t\cap(B_{p(t)})_{S_t}$. Comparing each copy with
that point yields
\[
 d_t\le4h+\theta(\rho_t+\rho_{p(t)}+\delta_t+\delta_{p(t)}).
\]
Use \eqref{decomp:bag-drift} for $t$ and its parent, then absorb the term
$\theta d_t$ using $\theta\le1/2$. With
$b_t=8h+2\theta(\rho_t+\rho_{p(t)})$, this gives
\begin{equation}\label{decomp:drift-recursion}
 d_t\le b_t+4\theta\sum_{j=1}^{k-1}d_{\operatorname{anc}_j(t)}.
\end{equation}

Write \eqref{decomp:drift-recursion} as $d\le b+\Gamma d$ entrywise, with
vectors indexed by nonroot bags.
The nonnegative matrix $\Gamma$ has nonzero entries only in ancestor
columns, so it is nilpotent. Its row sums are at most $4\theta(k-1)$ and its column
sums at most $4\theta D_k$: a node has at most $D_k$ descendants in the
next $k-1$ generations. For any matrix,
$\|\Gamma\|_2^2\le\|\Gamma\|_1\|\Gamma\|_\infty$; thus
\eqref{decomp:theta-conditions} gives $\|\Gamma\|_2\le1/2$.
Iterating the entrywise inequality and using nilpotence gives
$\|d\|_2\le2\|b\|_2$.
By $(a+b+c)^2\le3(a^2+b^2+c^2)$ and the child-degree bound,
\[
 \sum_tb_t^2\le12\left\{16Mh^2+
                \theta^2(1+\Delta)\sum_t\rho_t^2\right\}.
\]
Consequently
\begin{equation}\label{decomp:drift-squared}
 \sum_td_t^2\le48\{16Mh^2+
                       \theta^2(1+\Delta)kX_*^2\},
 \qquad
 \sum_t\delta_t^2\le K_1\sum_td_t^2.
\end{equation}
For the second assertion, square \eqref{decomp:bag-drift} by
Cauchy--Schwarz. An edge drift occurs for its own bag and for at most
$\sum_{j=1}^{k-2}\Delta^j$ descendants, giving precisely the bound $K_1$.

The chosen slopes have an exact telescoping identity:
\begin{equation}\label{decomp:slope-telescope}
 \sum_t\nabla a_t(x^*_{V_t})^{\mathsf T}(z^t-x_{V_t})
 =\sum_{t\ne r}\lambda_t^{\mathsf T}
                         (z^t_{S_t}-z^{p(t)}_{S_t}).
\end{equation}
To verify it, fix $i$ and telescope $z^t_i-x_i$ along its path to
$\operatorname{top}(i)$. The coefficient of an edge difference
$z^u_i-z^{p(u)}_i$ is the sum of $\partial_i a_t(x^*)$ over bags below
$u$ containing $i$. This is $\lambda_{u,i}$. Sum over $i$.
Thus the exact-factor version of a configuration value equals
\[
 f(x)+\sum_t\{a_t(z^t)-a_t(x_{V_t})-
               \nabla a_t(x^*_{V_t})^{\mathsf T}(z^t-x_{V_t})\}.
\]

Only at most $w$ separator coordinates differ in each bracket.
Lipschitz gradients, Taylor's formula on the segment between its two
points, and $\|z^t-x_{V_t}\|_2\le\sqrt w\,\delta_t$ bound the bracket
below by
$-H_a\sqrt w\,r_t\delta_t-(H_aw/2)\delta_t^2$.
The factor-relaxation loss in bag $t$ is at most
$\alpha'A\omega(B_t)^2/4$. Hence the relaxed configuration value is
at least $f(x)-E$, where
\[
 E\le\sum_t\left\{H_a\sqrt w\,r_t\delta_t+
      \frac{H_aw}{2}\delta_t^2+\frac{\alpha'A}{4}\omega(B_t)^2\right\}.
\]
Young's inequality gives
\[
 H_a\sqrt w\,r_t\delta_t
 \le\frac{c_g}{4k}r_t^2+\frac{H_a^2wk}{c_g}\delta_t^2.
\]
Also $\omega(B_t)^2\le3(4h^2+\theta^2\rho_t^2+
\theta^2\delta_t^2)$. Using \eqref{decomp:bag-radii} and $\theta\le1/2$
therefore gives
\[
 E\le\frac{c_g}{4}X_*^2+Q\sum_t\delta_t^2+
      3\alpha'AMh^2+\frac34\alpha'A\theta^2kX_*^2.
\]
Now apply \eqref{decomp:drift-squared} and
\eqref{decomp:theta-conditions}:
\[
 E\le\frac{3c_g}{4}X_*^2+MH_0h^2.
\]
Quadratic growth shows that every configuration is at least
$f_*+(c_g/4)X_*^2-MH_0h^2\ge f_*-\epsilon/4$.
Lemma~\ref{decomp:unfolding} proves the required root bound.
Finally, each bag partition has at most
$(J+1)(4/\theta)^{w+1}$ leaves and each separator partition no more than
that many cells. Summing gives \eqref{decomp:size-bound}.
\end{proof}

\subsection{The path separation}
\label{app:decomp:path}

\begin{proof}[Proof of Theorem~\ref{decomp:path-separation}]
The unary function $\varphi(t)=t^2-t^4/10$ has
$\varphi''(t)=2-6t^2/5\in[4/5,2]$ on $[-1,1]$.
Using $x_i^4\le x_i^2$ and
$2|x_ix_{i+1}|\le x_i^2+x_{i+1}^2$ gives
\begin{equation}\label{decomp:path-growth}
 f_n(x)\ge\tfrac1{10}\|x\|_2^2.
\end{equation}
Thus zero is the unique global minimizer. Its Hessian is the positive
definite tridiagonal matrix
$H_n=\operatorname{tridiag}(4/5,2,4/5)$, whose eigenvalues are
$2+(8/5)\cos(j\pi/(n+1))$, $1\le j\le n$.
At the all-ones point the Hessian instead has diagonal $4/5$ and the same
off-diagonals. Its smallest eigenvalue is
$4/5-(8/5)\cos(\pi/(n+1))<0$ for $n\ge3$. Continuity of the Hessian
gives a negative direction at nearby interior points, proving
nonconvexity. The bilinear alphaBB underestimator has Hessian
$\bigl(\begin{smallmatrix}4/5&4/5\\4/5&4/5\end{smallmatrix}\bigr)$,
so it is convex.

We first prove the accuracy-dependent spatial bound. Put
$\alpha_1=\alpha_n=2/5$ and $\alpha_i=4/5$ for $2\le i\le n-1$.
A valid box $B$ satisfies, at every $x\in B$,
\begin{equation}\label{decomp:path-validity}
 f_n(x)+\epsilon\ge\sum_i\alpha_i(x_i-l_i)(u_i-x_i).
\end{equation}
For a nondegenerate box, the arithmetic--geometric mean inequality
implies
\[
 \int_B(f_n+\epsilon)^{-n/2}\,dx
 \le n^{-n/2}\prod_i\alpha_i^{-1/2}
       \prod_i\int_{l_i}^{u_i}\frac{dx_i}
                    {\sqrt{(x_i-l_i)(u_i-x_i)}}
 =n^{-n/2}\pi^n\prod_i\alpha_i^{-1/2}.
\]
Degenerate boxes have zero $n$-dimensional measure. Summing over any
finite valid cover gives
\begin{equation}\label{decomp:weighted-integral}
 N\ge(n/\pi^2)^{n/2}\Bigl(\prod_i\alpha_i\Bigr)^{1/2}
                         \int_D(f_n+\epsilon)^{-n/2}\,dx.
\end{equation}

Since $f_n(x)\le\tfrac12x^{\mathsf T}H_nx$ and
$\lambda_{\min}(H_n)\ge2/5$, the ellipsoid
$\{x:\tfrac12x^{\mathsf T}H_nx\le1/5\}$ lies in the unit Euclidean ball
and hence in $D$. The substitution $x=\sqrt2H_n^{-1/2}y$ and polar
coordinates give from \eqref{decomp:weighted-integral}
\begin{equation}\label{decomp:path-exact-integral}
 N\ge2(2n/\pi)^{n/2}
       \left(\frac{\prod_i\alpha_i}{\det H_n}\right)^{1/2}
       \frac{J_n(\sqrt{1/(5\epsilon)})}{\Gamma(n/2)},
 \qquad
 J_n(T)=\int_0^Tt^{n-1}(1+t^2)^{-n/2}\,dt.
\end{equation}
The determinant recurrence
$\det H_j=2\det H_{j-1}-(16/25)\det H_{j-2}$, with initial values
$1,2$, yields
\[
 \det H_n=\tfrac43(8/5)^n-\tfrac13(2/5)^n
       \le\tfrac43(8/5)^n,
 \qquad \prod_i\alpha_i=(4/5)^n/4.
\]
For $T\ge\sqrt n$,
\[
 J_n(T)\ge e^{-1/2}\int_{\sqrt n}^T\frac{dt}{t}
         =\tfrac12e^{-1/2}\log(T^2/n),
\]
because $(1+t^{-2})^{-n/2}\ge e^{-n/(2t^2)}\ge e^{-1/2}$.
The standard upper Stirling bound
\[
 \Gamma(n/2)\le\sqrt{4\pi/n}\,(n/(2e))^{n/2}e^{1/(6n)}
\]
then bounds the coefficient of
$\sqrt n(2e/\pi)^{n/2}\log(1/(5n\epsilon))$ below by
$\sqrt3e^{-1/2-1/(6n)}/(8\sqrt\pi)>0.068$ for $n\ge3$.
This proves the first expression in \eqref{decomp:path-lower} when its
logarithm is nonnegative; otherwise that expression is below zero.

For the second bound, define
$d_i^B(x)=\min\{x_i-l_i,u_i-x_i\}$.
The alphaBB gap of each edge is at least
\[
 \tfrac25\{(x_i-l_i)(u_i-x_i)+(x_{i+1}-l_{i+1})(u_{i+1}-x_{i+1})\}
 \ge\tfrac45 d_i^B(x)d_{i+1}^B(x).
\]
Hence every valid box satisfies
$\tfrac45\sum_{i<n}d_i^B(x)d_{i+1}^B(x)\le f_n(x)+\epsilon$.
Intersect such a box with $D$, writing the intersection as
$\prod_i[z_i-h_i,z_i+h_i]$, with $0\le h_i\le1$.
If it has positive volume, put $s_i=1-h_i\in[0,1)$.
Its center obeys $|z_i|\le s_i$ and $d_i^B(z)\ge h_i$.
The upper quadratic bound on $f_n$ therefore implies
\[
 \tfrac45\sum_{i<n}(1-s_i)(1-s_{i+1})
 \le\sum_i s_i^2+\tfrac45\sum_{i<n}s_is_{i+1}+\epsilon.
\]
Cancel the cross terms and use that each index has degree at most two:
\begin{equation}\label{decomp:center-constraint}
 n-1\le\sum_i(\tfrac54s_i^2+2s_i)+\tfrac54\epsilon.
\end{equation}
For $0\le s<1$,
\begin{equation}\label{decomp:center-scalar}
 \log(1-s)\le\log(3/5)
          -\tfrac59(\tfrac54s^2+2s-1).
\end{equation}
Indeed, the derivative of
$\log(1-s)+(5/9)(5s^2/4+2s)$ is
$(2-5s)(1+5s)/(18(1-s))$, so its maximum is at $s=2/5$;
there $5s^2/4+2s=1$. Summing
\eqref{decomp:center-scalar} and using
\eqref{decomp:center-constraint} gives
\[
 \frac{\vol(B\cap D)}{\vol(D)}=\prod_i(1-s_i)
 \le(3/5)^n\exp[\tfrac59(1+\tfrac54\epsilon)].
\]
Zero-volume intersections obey the same bound. Volume subadditivity over
the cover proves the second expression in \eqref{decomp:path-lower}.

Both lower bounds also hold for rectangular owned certification events
$(A_e,B_e)$ covering $D$, with $A_e\subseteq B_e$ and the corresponding
gap test valid at owned points. For the integral proof, integrate over
$A_e$ and bound its arcsine integral by the integral over $B_e$.
For the volume proof, use the center of the closure of $A_e\cap D$.
If this owner has positive volume, its center is in its interior and
hence in the owner; distances to the faces of $B_e$ are at least its
half-widths. Open or closed endpoint choices do not change its volume.
Thus the same per-event estimates apply. This is the owned event count,
which can include removed-region certificates and virtual probe events.

For the decomposition upper bound, assign the unary factor at coordinate
$t$ and the edge $(t,t+1)$ to bag $t$; assign the final unary factor also
to the final bag. Then $M=n-1$, $w=1$, $k=2$, $\Delta=1$,
$D_k=K_1=1$, $c_g=1/10$, $H_a=14/5$, and $\alpha'A=4/5$.
The Hessian norm of each bag objective is at most $2+4/5=14/5$.
Thus $Q=159$ and $H_0=122114.4$. The choice $\theta=2^{-10}$ satisfies
\eqref{decomp:theta-conditions}. Choose
$h^2=\epsilon/(488457.6(n-1))$.
All slopes vanish because every factor gradient vanishes at zero.
Theorem~\ref{decomp:small-certificate} gives
\[
 N_{\mathrm{dec}}\le2(n-1)4096^2
 \left\{\tfrac12\log_2\frac{1953830.4(n-1)}{\epsilon}+2\right\},
\]
which is bounded by \eqref{decomp:path-upper}.

For the uniform ratio, let $a=\sqrt{2e/\pi}$ and split into two cases.
If $\epsilon\le0.04/n^2$, then
$\log(0.2/(n\epsilon))\ge\tfrac12\log(1/\epsilon)$ and
$n\le0.2\epsilon^{-1/2}$. Since $\epsilon\le10^{-4}$,
\eqref{decomp:path-upper} is at most $C_1n\log(1/\epsilon)$ for an
absolute constant $C_1$. The first spatial bound consequently gives
$N/U_n\ge c_1a^n/\sqrt n$.
If $0.04/n^2<\epsilon\le10^{-4}$, then $1/\epsilon<25n^2$,
so $U_n\le C_2n\log n$. The second spatial bound gives
$N/U_n\ge c_2(5/3)^n/(n\log n)$.
Since $(5/3)/a>1$, the positive sequence
$((5/3)/a)^n/(\sqrt n\log n)$ tends to infinity. Its infimum over
integers $n\ge3$ is therefore positive. Decreasing the constant proves
\eqref{decomp:uniform-ratio} in both cases.

Finally, for $b>0$ the McCormick gap on a rectangle equals
\[
 b\min\{(x-l_x)(y-l_y),(u_x-x)(u_y-y)\}.
\]
It is at least $b\,d_xd_y$, which is all the center-volume proof needs.
It is at most
$\tfrac b2\{(x-l_x)(u_x-x)+(y-l_y)(u_y-y)\}$:
use $\min\{a_1b_1,a_2b_2\}\le\sqrt{a_1a_2b_1b_2}
\le(a_1a_2+b_1b_2)/2$ with
$a_1=x-l_x$, $a_2=u_x-x$, $b_1=y-l_y$, $b_2=u_y-y$.
Thus it satisfies the upper error hypothesis with $\alpha'=2/5$,
and the decomposition upper bound is unchanged. The strictly positive
coordinate gap in \eqref{decomp:path-validity} is not available, so the
accuracy-dependent integral bound is not transferred.
\end{proof}
